\subsection{Residual systematic effects} %% the biases that we observe in chapter 5
\label{lab:Syst_ResidualEffecte}

%* 1. Discuss scambled toy studies.
%   Why we expect the average fit result to be zero,
%   indead we observed it consistent with zero in toy sutdies.
In Section~\ref{app:toys_signal_injected_scrambled} we tested our scrambled method on
our toy samples. The toys consist of SM-like toys (No signals), and 
signal-injected toys. In sidereal phase, the MS-like toys show data points
are distributed according to only statistical fluctuation. The signal injected toys show 
some oscillations correspond to injected signal shape.

The figure~\ref{fig:SM_like_toy_scrambled} shows a comparison of a SM-like toy before (left) and
and after scrambled (right). We can see that the scrambled toy is just like an original one,
it's also SM-like toy. In the same way we tested the scrambling on 1000 SM-like toys. 
The figure~\ref{fig:SM_like_toy_fit_results12} shows average of 1000 fitted results and their one/two 
sigma standard deviation. The left plot shows the results from SM-like toys before scrambling,
and the right plot shows the results after apply scrambling.
As we expected, in both case, average fit result on each coefficeint is in a good 
agreement with zero. Because, in SM-like toy datasets we have only 
simulated statistical uncertainties. Therefore, we expect no bias 
(average fit result consistent with zero).

In signal injected case, scrambling can completely washed the injected singal. 
The figure~\ref{fig:signal_injected_toy_p003_beforeAfter_scramble} show the compareson 
of before and after scrambling of an siganl injected toy sample. In the left plot we clearly 
see the oscillation pateren, but in the right plot after scrambling we didn't observe presence 
of any signals. After scrambing the fitted signal strength down to $-4.4E^{-7}\pm 3.8E^{-6}$;
 it was $5.3E^{-5}\pm 3.8E^{-6}$ before scrambling. This study concluded that event 
 if a singal is present the scrambling can wash it.
Also, if a dataset has only statistical uncertaity (e.g other systematics effects are negligible),
we can still expect that the fitted signal strength after scrambling will be 
consistent with zere (e.g no bias will be seen).



%* 2. Discuss the scrambled results from partially unblinded data.
%   We have observed the average fit results in several coefficients
%   diviate from zero at the order of one(or less) STD level. This is
%   different than what we learned from toy studies.
However, when we check the data using scrambling (partially unblind), we observe 
that average values of fit results in some singal strengths (coefficents) deviate 
from zero (bias). This can be clearly seen in figure~\ref{fig:real_data_scrambled_summary_combined}
where Run 2 - combined fit results are presented. In the plot, each black dot is an average fit results 
of 1000 scrabled datasets. In a few coefficents we observed the average fit results deviate from zero
but still within one sigma range. The exact numbers can be read from the table~\ref{table:scrambel_results_ll}.
This kind of deviation indicate that there is a systematic effact other than statisticle fluctuation. 
Because, as we expected from our toy studies we discussed before, if the statistical fluctuation is the 
only source of uncertainty then we should observe results from scrambled dataset also consistent with zero.

%* 3. The difference between the average fit result and zero (zero is
%    the theory expectation) in each coefficient can be a measurement of
%    spurious signal uncertaity. 
The deviation from zero in each coefficeint shown in figure~\ref{fig:real_data_scrambled_summary_combined} can be seen
as size of spurious signal uncertainty for corresponding coefficent.
 The spurious signal is models as Gaussian distribution centered at zero and it standard deviation
  is absolute values of average fit results (size of the deviation from zero) shown in figure~\ref{fig:real_data_scrambled_summary_combined}.
  The size of spurious signals for Run 2 fit and Run 2 - combined fit are given in table~\ref{table:spurious_signal_ll}.
  

% Sys table
%%%%%%%%%%%%%%%
%%ee+uu channel%%%
%%%%%%%%%%%%%%%
\begin{table}[hp]
  \begin{center}
  %%\resizebox{\textwidth}{!}{
  \begin{tabular}{|l|c| c| c| c |c|}
  \hline
   & Run2 & Run2-Comb. Fit \\
  \hline
  $d_{\it u}^{\it X,Z}$   & -3.23e-06 & -2.99e-06  \\
  \hline
  $d_{\it u}^{\it Y,Z}$   & -1.69e-06  & -2.5e-06  \\
  \hline
  $d_{\it u}^{\it XY,XY}$ & 8.03e-07  & -2.39e-07  \\
  \hline
  $d_{\it u}^{\it X,Y}$  & 3.21e-07  & 1.41e-07  \\
  \hline
  $c_{\it u}^{\it X,Z}$  & -2.51e-06  & -2.32e-06 \\
  \hline
  $c_{\it u}^{\it Y,Z}$  & -1.32e-06  & -1.94e-06  \\
  \hline
  $c_{\it u}^{\it XY,XY}$ & 6.24e-07  & -1.86e-07  \\
  \hline
  $c_{\it u}^{\it X,Y}$  & 2.50e-07  & 1.09e-07 \\
  \hline
  $c_{\it d}^{\it X,Z}$  & -1.12e-04  & -1.03e-04 \\
  \hline
  $c_{\it d}^{\it Y,Z}$  & -5.86e-05 & -8.66e-05  \\
  \hline
  $c_{\it d}^{\it XY,XY}$  & 2.78e-05 & -8.27e-06 \\
  \hline
  $c_{\it d}^{\it X,Y}$ & 1.11e-05 & 4.86e-06 \\
  \hline
  \hline
  \end{tabular}
  \end{center}
  \caption{
  Size of spurious signal uncertainties in $ee+\mu\mu$ channel, for Run 2 fit (first column) 
  and Run 2 - combined fit (second column, obtained from 10K samples).
  }\label{table:spurious_signal_ll}
  \end{table}

%%% Copy profiled X2 fit to here?





